\documentclass[a4paper,11pt]{article}
\usepackage[T2A]{fontenc}
\usepackage[utf8]{inputenc}
\usepackage[russian]{babel}
\usepackage{amsmath,amssymb,mathtools}
\usepackage{geometry}
\geometry{left=23mm,right=23mm,top=22mm,bottom=23mm}
\usepackage{microtype}
\usepackage{icomma}
\usepackage{tikz,tkz-euclide}
\usetikzlibrary{arrows.meta,positioning,shapes.geometric,calc}
\usepackage{tabularx,booktabs,array,enumitem}
\usepackage{xcolor}
\usepackage{hyperref}
\usepackage{fancyhdr}
\usepackage{parskip}
\usepackage{needspace}
\definecolor{accent}{HTML}{087E83}
\definecolor{darkaccent}{HTML}{07545A}
\definecolor{ink}{HTML}{292929}
\definecolor{faint}{HTML}{637579}
\hypersetup{colorlinks=true,linkcolor=darkaccent,pdfauthor={Лёвин Артём Александрович},pdftitle={Чек-лист: призма — Ксения Клыкова}}
\color{ink}
\setlength{\parindent}{0pt}
\setlength{\parskip}{5.3pt}
\renewcommand{\arraystretch}{1.25}
\setlist[itemize]{leftmargin=1.6em,itemsep=2pt,topsep=3pt}
\setlist[enumerate]{leftmargin=1.7em,itemsep=3pt,topsep=3pt}
\pagestyle{fancy}
\fancyhf{}
\fancyhead[L]{\small\color{faint}Призма: объём и поверхность}
\fancyhead[R]{\small\color{faint}Ксения Клыкова}
\fancyfoot[C]{\small\color{faint}\thepage}
\renewcommand{\headrulewidth}{0.28pt}
\renewcommand{\headrule}{\hbox to\headwidth{\color{accent!60}\leaders\hrule height \headrulewidth\hfill}}
\setlength{\headheight}{23pt}
\newcommand{\sbase}{S_{\text{осн}}}
\newcommand{\slat}{S_{\text{бок}}}
\newcommand{\stot}{S_{\text{полн}}}
\newcommand{\per}{P_{\text{осн}}}
\newcommand{\checkitem}{\(\square\)}
\newcommand{\key}[1]{\textcolor{darkaccent}{\textbf{#1}}}
\newcommand{\sechead}[1]{\Needspace{9\baselineskip}\vspace{7pt}{\Large\color{darkaccent}\bfseries #1}\par\vspace{3pt}}
\newcommand{\subhead}[1]{\Needspace{6\baselineskip}\vspace{5pt}{\large\color{darkaccent}\bfseries #1}\par}
\begin{document}
\thispagestyle{empty}
\begin{center}
\vspace*{24mm}
{\Large\color{darkaccent}МАТЕМАТИКА}\par
\vspace{20mm}
{\Huge Чек-лист}\par
\vspace{11mm}
{\LARGE\color{accent}Призма: объём, площадь поверхности\\[3pt] и изменение размеров}\par
\vspace{17mm}
{\large Подготовка к ЕГЭ по профильной математике}\par
\vfill
{\Large Ксения Клыкова}\par
\vspace{12mm}
{\large 23 ноября 2026 г.}\par
\vspace{18mm}
{\normalsize Лёвин Артём Александрович\\[3pt]эксклюзивно для Ксении Клыковой}\par
\vspace{15mm}
\begin{tikzpicture}[x=1cm,y=1cm]
 \draw[accent,line width=0.7pt] (0,0) .. controls (1.25,0.57) and (2.5,-0.57) .. (3.75,0)
 .. controls (5,0.57) and (6.2,-0.57) .. (7.5,0);
 \fill[darkaccent] (3.75,0) circle (1.2pt);
\end{tikzpicture}
\end{center}
\clearpage
\sechead{1. Что называется призмой}
\textbf{Призма} --- многогранник, у которого два основания представляют собой равные многоугольники в параллельных плоскостях, а боковые грани являются параллелограммами.
\begin{itemize}
 \item \key{Основания} могут быть треугольниками, четырёхугольниками, шестиугольниками и другими многоугольниками.
 \item \key{Боковые рёбра} соединяют соответствующие вершины оснований.
 \item \key{Высота \(H\)} --- расстояние между плоскостями оснований; её измеряют по перпендикуляру к ним.
 \item \key{Прямая призма:} боковые рёбра перпендикулярны основаниям, поэтому \(H=l_{\text{бок. ребра}}\). Её боковые грани --- прямоугольники.
 \item \key{Наклонная призма:} боковое ребро обычно длиннее высоты. Формула \(\slat=\per H\) для неё в общем случае неприменима.
\end{itemize}
\begin{center}
\begin{tikzpicture}[scale=.89,line join=round]
 \coordinate (A) at (0,0);\coordinate (B) at (3.4,0);\coordinate (C) at (1.05,1.05);
 \coordinate (AA) at (0,2.55);\coordinate (BB) at (3.4,2.55);\coordinate (CC) at (1.05,3.60);
 \draw[accent,very thick] (A)--(B)--(C)--cycle (AA)--(BB)--(CC)--cycle;
 \draw[accent,very thick] (A)--(AA) (B)--(BB) (C)--(CC);
 \draw[darkaccent,line width=1.5pt] (A)--node[left]{\(H\)}(AA);
 \fill (A) circle (1.2pt) node[below left]{\(A\)};
 \fill (B) circle (1.2pt) node[below right]{\(B\)};
 \fill (C) circle (1.2pt) node[left]{\(C\)};
 \fill (AA) circle (1.2pt) node[above left]{\(A_1\)};
 \fill (BB) circle (1.2pt) node[right]{\(B_1\)};
 \fill (CC) circle (1.2pt) node[above]{\(C_1\)};
 \node[align=left,text width=5.6cm,anchor=west] at (5.2,1.95) {\textbf{Прямая призма}\\[4pt]Основания: \(ABC\) и \(A_1B_1C_1\).\\Боковые грани: \(ABB_1A_1\), \(BCC_1B_1\), \(ACC_1A_1\).};
\end{tikzpicture}
\end{center}
\textbf{Правильная треугольная призма} --- прямая призма с \emph{равносторонним} треугольником в основании. Более общо: правильная призма прямая и имеет правильный многоугольник в основании.
\sechead{2. Три главные формулы}
\[
 \boxed{V=\sbase H} \qquad
 \boxed{\slat=\per H}\quad\text{(для прямой призмы)}
\]
\[
 \boxed{\stot=\slat+2\sbase}
\]
\textcolor{faint}{Здесь \(V\) --- объём; \(\sbase\) --- площадь одного основания; \(\per\) --- его периметр; \(\slat\) --- сумма площадей всех боковых граней.}
\par\medskip
\sechead{3. Площади оснований: формулы, которые нужно знать}
\begin{center}
\begin{tabularx}{\linewidth}{@{}>{\raggedright\arraybackslash}p{0.40\linewidth}>{\centering\arraybackslash}X@{}}
\toprule
Фигура в основании & Площадь основания \\
\midrule
Равносторонний треугольник со стороной \(a\) & \(\displaystyle \sbase=\frac{\sqrt3}{4}a^2\) \\
Правильный шестиугольник со стороной \(a\) & \(\displaystyle \sbase=\frac{3\sqrt3}{2}a^2\) \\
Ромб с диагоналями \(d_1,d_2\) & \(\displaystyle \sbase=\frac{d_1d_2}{2}\) \\
Квадрат со стороной \(a\) & \(\displaystyle \sbase=a^2\) \\
Прямоугольный треугольник с катетами \(a,b\) & \(\displaystyle \sbase=\frac{ab}{2}\) \\
\bottomrule
\end{tabularx}
\end{center}
\subhead{Пример 1. Объём правильной треугольной призмы}
\textbf{Условие.} Все рёбра правильной треугольной призмы равны \(\sqrt3\). Найти объём.

Поскольку призма прямая и все рёбра равны, \(a=H=\sqrt3\). Подставляем в формулу объёма:
\[
 V=\sbase H=\frac{\sqrt3}{4}a^2 H
 =\frac{\sqrt3}{4}(\sqrt3)^2\sqrt3
 =\frac{\sqrt3\cdot3\cdot\sqrt3}{4}
 =\boxed{\frac94}.
\]
\textbf{Проверка:} \((\sqrt3)^2=3\), а \(\sqrt3\cdot\sqrt3=3\).
\subhead{Ромб: как найти сторону по диагоналям}
У ромба диагонали \emph{перпендикулярны и делятся точкой пересечения пополам}. Они не обязаны быть равными.
\begin{center}
\begin{tikzpicture}[scale=.82]
\tkzDefPoint(-2,0){A}\tkzDefPoint(0,1.5){B}\tkzDefPoint(2,0){C}\tkzDefPoint(0,-1.5){D}\tkzDefPoint(0,0){O}
\tkzDrawPolygon[color=accent,thick](A,B,C,D)
\tkzDrawSegments[dashed,color=faint](A,C B,D)
\tkzDrawPoints[fill=darkaccent,size=2](A,B,C,D,O)
\tkzLabelPoints[left](A)\tkzLabelPoints[right](C)\tkzLabelPoints[above](B)\tkzLabelPoints[below](D)\tkzLabelPoints[above right](O)
\tkzMarkRightAngle[size=.25,color=darkaccent](B,O,C)
\node[anchor=west,align=left] at (3.0,.2) {\(AO=\dfrac{AC}{2}\), \(BO=\dfrac{BD}{2}\)\\[9pt]
 \(AB=\sqrt{AO^2+BO^2}\)\\[9pt]
 \(\sbase=\dfrac{AC\cdot BD}{2}\)};
\end{tikzpicture}
\end{center}
Если диагонали равны 4 и 3, то половины диагоналей --- 2 и \(1{,}5\), поэтому
\[
 a=\sqrt{2^2+1{,}5^2}=\sqrt{6{,}25}=2{,}5,\qquad \per=4a=10.
\]
\par\medskip
\sechead{4. Полная и боковая поверхность}
Для \emph{прямой} призмы площади боковых прямоугольников равны произведениям соответствующих сторон основания на высоту. Поэтому
\[
 \slat=(a_1+a_2+\cdots+a_n)H=\per H.
\]
Для полной поверхности к ним добавляем \textbf{два} равных основания:
\[
 \stot=\per H+2\sbase.
\]
\subhead{Пример 2. Ромб в основании}
\textbf{Условие.} Основание прямой призмы --- ромб с диагоналями 4 и 3. Боковое ребро равно 5. Найти площадь полной поверхности.
\begin{enumerate}
 \item По свойству ромба половины диагоналей равны 2 и \(1{,}5\).
 \item Сторона ромба: \(a=\sqrt{2^2+1{,}5^2}=2{,}5\).
 \item Периметр основания: \(\per=4\cdot2{,}5=10\).
 \item Поскольку призма прямая, \(H=5\). Следовательно, \(\slat=10\cdot5=50\).
 \item Площадь одного основания: \(\sbase=\frac12\cdot4\cdot3=6\).
 \item Полная поверхность:
 \[
 \stot=50+2\cdot6=\boxed{62}.
 \]
\end{enumerate}
\subhead{Пример 3. Найти высоту через полную поверхность}
\textbf{Условие.} Основание прямой призмы --- квадрат со стороной 3; полная поверхность равна 66. Найти высоту.
\[
 \sbase=3^2=9,\qquad \per=4\cdot3=12.
\]
\[
 66=12H+2\cdot9\quad\Longrightarrow\quad
 12H=66-18=48\quad\Longrightarrow\quad\boxed{H=4}.
\]
\textbf{Проверка:} \(12\cdot4+2\cdot9=48+18=66\).
\sechead{5. Как читать условие задачи}
\begin{itemize}
 \item «Найдите \textbf{боковую} поверхность» --- считаем \(\slat\): все боковые грани, без оснований.
 \item «Найдите \textbf{полную} поверхность» или просто «площадь поверхности» --- считаем \(\stot\): боковые грани и оба основания.
 \item «Найдите объём» --- считаем \(V=\sbase H\).
\end{itemize}
\par\medskip
\sechead{6. Отсечённая призма и средняя линия}
Пусть \(ABC\) --- треугольное основание прямой призмы, \(M\) --- середина \(AB\), \(T\) --- середина \(AC\), а через \(MT\) проведена плоскость, параллельная боковым рёбрам. Получена меньшая прямая призма с основанием \(AMT\).
\begin{center}
\begin{tikzpicture}[scale=1.0]
\tkzDefPoint(0,0){A}\tkzDefPoint(4.4,0){B}\tkzDefPoint(1.45,2.55){C}
\tkzDefMidPoint(A,B)\tkzGetPoint{M}
\tkzDefMidPoint(A,C)\tkzGetPoint{T}
\tkzDrawPolygon[color=accent,thick](A,B,C)
\tkzDrawSegments[color=darkaccent,very thick](M,T)
\tkzDrawPoints[fill=darkaccent,size=2](A,B,C,M,T)
\tkzLabelPoints[below](A,B,M)\tkzLabelPoints[above](C)\tkzLabelPoints[left](T)
\node[anchor=west,align=left,text width=6.2cm] at (5.2,1.6) {\(AM=\dfrac{AB}{2}\),\\[8pt]
\(AT=\dfrac{AC}{2}\),\\[8pt]
\(MT=\dfrac{BC}{2}\).};
\end{tikzpicture}
\end{center}
\textbf{Ключевой факт:} средняя линия треугольника параллельна третьей стороне и равна её половине. Поэтому периметр малого треугольника в 2 раза меньше:
\[
 P_{AMT}=AM+AT+MT=\frac{AB+AC+BC}{2}=\frac{P_{ABC}}{2}.
\]
\subhead{Пример 4. Как меняется боковая поверхность}
Высоты малой и исходной призм одинаковы, обозначим их \(H\). Для двух призм пишем одинаковую формулу:
\[
 \slat^{(\text{мал})}=P_{AMT}H,\qquad
 \slat^{(\text{бол})}=P_{ABC}H.
\]
Делим первую площадь на вторую, сокращаем \(H\):
\[
 \frac{\slat^{(\text{мал})}}{\slat^{(\text{бол})}}
 =\frac{P_{AMT}H}{P_{ABC}H}
 =\frac{P_{ABC}/2}{P_{ABC}}
 =\frac12.
\]
Если \(\slat^{(\text{бол})}=12\), то
\[
 \boxed{\slat^{(\text{мал})}=\frac12\cdot12=6}.
\]
\textbf{Важно:} коэффициенты для различных величин отличаются. Стороны основания уменьшились в 2 раза, его площадь --- в 4 раза, боковая поверхность при неизменной высоте --- в 2 раза.
\par\medskip
\sechead{7. Изменение размеров призмы: объём}
Если все \emph{линейные размеры основания} увеличены в \(k\) раз, его площадь увеличивается в \(k^2\) раз. Если при этом высота изменилась в \(m\) раз, то объём изменится в \(k^2m\) раз:
\[
 \boxed{\frac{V_2}{V_1}=k^2m},\qquad
 k=\frac{a_2}{a_1},\quad m=\frac{H_2}{H_1}.
\]
Формула применима, если новые основания подобны старым; \(k>0\), \(m>0\).
\subhead{Пример 5. Сторона втрое больше, высота вдвое меньше}
\textbf{Условие.} Объём правильной треугольной призмы равен 6. Как изменится объём, если сторону основания увеличить в 3 раза, а высоту уменьшить в 2 раза?
\[
 a_2=3a_1,\qquad H_2=\frac{H_1}{2}.
\]
Так как площадь равностороннего треугольника равна \(\frac{\sqrt3}{4}a^2\), имеем
\[
 V_1=\frac{\sqrt3}{4}a_1^2H_1,\qquad
 V_2=\frac{\sqrt3}{4}a_2^2H_2.
\]
Делим вторую формулу на первую и подставляем изменения:
\[
 \frac{V_2}{V_1}=
 \frac{a_2^2H_2}{a_1^2H_1}
 =\frac{(3a_1)^2(H_1/2)}{a_1^2H_1}
 =\frac{9}{2}.
\]
Поэтому
\[
 \boxed{V_2=6\cdot\frac92=27}.
\]
\textbf{Два независимых эффекта:} площадь основания увеличилась в 9 раз, высота уменьшилась в 2 раза. Совместно это даёт увеличение объёма в \(9/2=4{,}5\) раза.
\subhead{Общий способ: метод отношений}
\begin{enumerate}
 \item Запишите формулы искомой величины \emph{до} и \emph{после} изменения.
 \item Разделите вторую формулу на первую.
 \item Подставьте отношения соответствующих размеров.
 \item Сократите общие множители, найдите отношение величин.
 \item Умножьте известное исходное значение на полученный коэффициент.
\end{enumerate}
\textcolor{faint}{Этот приём работает при сравнении объёмов, площадей и других однотипных величин; важно правильно определить, что именно менялось.}
\clearpage
\thispagestyle{fancy}
\begin{center}
{\Large\color{darkaccent}\bfseries Алгоритм: сравнение величин}\par\vspace{10mm}
\begin{tikzpicture}[
 node distance=12mm,
 act/.style={draw=accent,rounded corners=3pt,minimum height=11mm,text width=11.2cm,align=center,font=\normalsize,inner sep=7pt},
 term/.style={draw=darkaccent,ellipse,minimum height=10mm,text width=6.3cm,align=center,font=\normalsize,inner sep=5pt},
 arr/.style={-{Stealth[length=3.0mm]},darkaccent,line width=.7pt}
]
\node[term] (start) {Есть исходная и изменённая фигуры};
\node[act,below=of start] (a) {Выделите известную величину и запишите\quad \(Q_1=F(\text{размеры}_1)\)};
\node[act,below=of a] (b) {Запишите ту же формулу для новой фигуры\quad \(Q_2=F(\text{размеры}_2)\)};
\node[act,below=of b] (c) {Составьте отношение \(\displaystyle\frac{Q_2}{Q_1}\)};
\node[act,below=of c] (d) {Подставьте коэффициенты изменения размеров и сократите общие множители};
\node[act,below=of d] (e) {Получите \(\displaystyle\frac{Q_2}{Q_1}=r\), затем \(Q_2=rQ_1\)};
\node[term,below=of e] (finish) {Проверьте смысл и единицы измерения};
\draw[arr] (start)--(a);\draw[arr] (a)--(b);\draw[arr] (b)--(c);\draw[arr] (c)--(d);\draw[arr] (d)--(e);\draw[arr] (e)--(finish);
\end{tikzpicture}
\end{center}
\clearpage
\sechead{8. Типичные ошибки и как их избежать}
\begin{tabularx}{\linewidth}{@{}>{\raggedright\arraybackslash}p{0.40\linewidth}X@{}}
\toprule
Ошибка & Правильная проверка \\
\midrule
Путать площадь основания и периметр & \(\sbase\) измеряется в квадратных единицах; \(\per\) --- в линейных. \\
Использовать \(\slat=\per H\) для любой призмы & Убедитесь, что призма \emph{прямая}. \\
Забывать второе основание & В полной поверхности присутствует \(2\sbase\). \\
Путать площадь квадрата и его сторону & При стороне 3 площадь \(3^2=9\); периметр \(4\cdot3=12\). \\
Терять плюс в формуле полной поверхности & \(\stot=\slat+2\sbase\); проверьте знак отдельно. \\
При уменьшении высоты вдвое писать \(H_2=2H_1\) & Верно: \(H_2=H_1/2\), или \(H_1=2H_2\). \\
Увеличивать площадь основания втрое при увеличении стороны втрое & При подобии площадей множитель возводится в квадрат: \(3^2=9\). \\
\bottomrule
\end{tabularx}
\sechead{9. Итоговый чек-лист перед решением}
\begin{itemize}
 \item[\checkitem] Я определил(а) вид призмы: прямая или наклонная.
 \item[\checkitem] Я указал(а), какую площадь надо найти: основание, боковую или полную поверхность.
 \item[\checkitem] Я выбрал(а) формулу площади основания и правильно подставил(а) размеры.
 \item[\checkitem] Я помню, что объём равен \(\sbase H\), а для прямой призмы \(\slat=\per H\).
 \item[\checkitem] Я проверил(а) коэффициенты масштабирования: стороны, площади, объём.
 \item[\checkitem] Я отдельно перепроверил(а) знаки, степени, единицы измерения и арифметику.
\end{itemize}
\vspace{5pt}
\begin{center}
\begin{tikzpicture}
 \draw[accent,line width=.65pt] (0,0) .. controls (1.1,.23) and (2.1,-.23) .. (3.2,0);
\end{tikzpicture}\\[-2pt]
{\small\color{faint}Лёвин Артём Александрович · эксклюзивно для Ксении Клыковой}
\end{center}
\end{document}
